\subsection*{Spectral benchmark on $S^3$ and Hopf-compatible diagnostics}
  \label{subsec:hopf_spectral}

  The ratio $\lambda_2/\lambda_1=8/3$ is an exact spectral property of
  the scalar Laplace--Beltrami operator on the unit three-sphere.
  The Hopf fibration $S^1\hookrightarrow S^3\to S^2$ does not determine
  this ratio; it provides only a canonical fiber/base partition used
  to define the kernel-weighted observable $R(\alpha)$.

  \paragraph{Exact spectral benchmark.}
    For $\Delta_{S^3}$ with unit radius,
    \[
      \lambda_k = k(k+2), \qquad \text{mult. }(k+1)^2,
    \]
    so that $\lambda_1=3$, $\lambda_2=8$, and
    \[
      \frac{\lambda_2}{\lambda_1}=\frac{8}{3}.
    \]
    This ratio is intrinsic and independent of any fibration
    (Appendix~\ref{sec:analytical_ratio}).

  \paragraph{Hopf-compatible splitting.}
    The Hopf structure induces a geometric decomposition into
    fiber and base directions, motivating the separation
    $d_{\mathrm{fiber}}^2$ and $d_{\mathrm{base}}^2$
    in the anisotropic kernel $K_\alpha$.
    The associated $k=1$ eigenspace (dimension $4=(1+1)^2$)
    may, once a Hopf/base symmetry is fixed, be organized as
    a singlet plus a triplet under $\mathrm{SO}(3)$ acting on the base.
    This representation-theoretic decomposition is interpretative
    and does not affect the intrinsic eigenvalue ratio
    $\lambda_2/\lambda_1 = 8/3$.

  \paragraph{Connection to $R(\alpha)$.}
    The observable
    $R(\alpha)$ biases edge contributions toward fiber-sensitive
    increments as $\alpha$ increases.
    In a Hopf-compatible projectable regime,
    the normalized response converges to the benchmark $8/3$.
    The kernel does not generate this value; it isolates
    the corresponding spectral structure while suppressing
    non-projectable contributions.

  \paragraph{Conclusion.}
    The value $8/3$ is a topological spectral invariant of $S^3$.
    $R(\alpha)$ provides an operational diagnostic that approaches
    the same benchmark under refinement, without asserting an
    operator-level identity between the two quantities.
